% Condensed evaluation for the short edition. Every figure is copied from a
% retained record named beside it; the full narratives, the discarded
% attempts and the per-attempt tables are in the extended edition.
\section{Evaluation: the executed campaigns}
\label{sec:evaluation}

This section reports what the reference implementation~\cite{dispat} was
observed to do in the campaigns that were executed, with the binary, the
environment and the counting unit of each stated separately. Five questions
organise it. \evRQ{1} asks whether the consumer-relative admission window
retains the residual obligations that a source-relative window loses.
\evRQ{2} asks whether versions and completion records stay stable through
repeated failure, repair and identity conflict, and how far the plans agree
with an independent model. \evRQ{3} asks whether execution respects artefact
availability and CI gates. \evRQ{4} asks what planning costs as the graph,
the intent and the history grow. \evRQ{5} asks what the distributed execution
profile of \spec{28} measured on two real build graphs, and what it did not.
Every number is transcribed from a retained file that the sentence or table
note names; no campaign is pooled with another; designed deterministic
matrices are reported as counts with their denominators and receive no
$p$-values. The extended edition carries the full account of each campaign,
including every failed and discarded attempt.

\subsection{Evidence boundaries and units}
\label{sec:eval:boundaries}

\Cref{tab:eval:boundaries} pins each body of evidence to its binary,
environment and counting unit. The units are not interchangeable: a
\emph{history} is one generated repository planned once; an \emph{episode} is
one fixture driven through fault, retry and repair; a \emph{state} or
\emph{transition} is an enumeration unit of the finite model; a
\emph{predicate} is one verifier check applied repeatedly to fixtures; a
\emph{release} is one \code{dispat release} from invocation to exit under GNU
\code{time}. Predicate counts are implementation checks, not independent
observations, and are never summed across campaigns. Every campaign through
1.10.0 runs the CCME 2.0.0 parser; the kernel and LLVM rounds run the
\spec{28} profile of the 1.11.0 release candidate and are the only evidence
about it.

\begin{center}
\footnotesize
\begin{longtable}{@{}L{2.6cm}L{6.0cm}L{5.4cm}@{}}
\caption{Evidence boundaries. Sources:
\protect\evpath{experiments/README.md} and
\protect\evpath{experiments/results/compare-v140-v160.md} (historical);
\protect\evpath{output/experiments-v182-summary.json} (1.8.2);
\protect\evpath{paper/saner/evidence/EXPERIMENTS-1.9.0-20260909.md} and
\protect\evpath{EXPERIMENTS-1.10.0-20260909.md} (companions);
\protect\evpath{experiments/saner-2027/campaign-mechanisms-20260910/RESULTS.md}
(mechanisms); the forks' \protect\evpath{dispat/gce/PARITY-20260922*.md} and
\protect\evpath{.ci/dispat/gce/REPORT.md}, copied under
\protect\evpath{experiments/distributed-20260922/evidence/} (kernel, LLVM).}
\label{tab:eval:boundaries}\\
\toprule
Campaign & Binary under test and environment & Counting unit and record \\
\midrule
\endfirsthead
\toprule
Campaign & Binary under test and environment & Counting unit and record \\
\midrule
\endhead
Historical differential, from 30 August 2026 & Released \code{dispat 1.4.0}
from the official image (release commit \code{5e3a2337}); repeated on 1.6.0,
1.7.0 and 1.7.2; Docker on an Apple-silicon host & History;
\evpath{experiments/results/v140}, \evpath{v160}, \evpath{v170}, \evpath{v172} \\
1.8.2 E1 to E5, 7 September 2026 & \code{dispat 1.8.2}, CLI commit
\evhash{a7011cbcb9e88e52e7336f6fc8f1b207a8e7a311}; Docker, two CPUs and 4~GiB
per container, contended host & History (E1), repetition (E2), release row
(E3), episode and timing cell (E4), check (E5);
\evpath{experiments/results/v182-20260907}, 753 files under its manifest \\
1.9.0 and 1.10.0 companions, 9 September 2026 & \code{dispat 1.9.0},
SHA-256 \evhash{be3370c472cdbfe66fe4256d3fc1ff78ff19bdd5f041ef222187a52767fa4f35};
\code{dispat 1.10.0}, SHA-256
\evhash{03d1ba5fbaef35b69582e9489f40b5438f81ae6926e02ea6d02716b5b73d92fd};
Docker Desktop Linux ARM64, two CPUs, 4~GiB, private registries only &
History, episode, execution, CI repetition, predicate;
\evpath{experiments/saner-2027/campaign-next-release-prepared/runs} (one failed
and one completed attempt) and \evpath{campaign-1.10.0-20260909} (9,157
archived files) \\
Mechanism matrix, 9 September 2026 & The 1.10.0 binary and runner image;
network disabled & Episode, state, transition;
\evpath{experiments/saner-2027/campaign-mechanisms-20260910}, 16,464 fixture
files archived \\
Kernel farm rounds, 22 September 2026 & Round 1: a development build, SHA-256
\evhash{a7172a7feb21e27ce882050120b9e04df1e6d4b4ee5f3341ab8a99ba155ba70a};
round 2: the feature branch at \code{f0c9e959}, static linux/amd64, SHA-256
\evhash{63441cc66f68d07b5efbf9a42edd5559304d153d48b82b1dfb037349e20606ed};
Google Compute Engine~\cite{gce}, \code{c3-standard-22} and
\code{t2d-standard-8}, us-central1 & Release (GNU \code{time}), per-cell
script seconds; \evpath{dispat/gce} of the fork at \code{c699bf3d}, copied
under \evpath{evidence/kernel-gce} \\
LLVM pair, 22 September 2026 & Attempts 1 to 8: \code{f0c9e959}; attempt 9:
\code{1.11.0-rc.4+main.272772cb} as printed in its release log; Compute
Engine, \code{e2-standard-4} orchestrator, two \code{n2-custom-14-49152}
workers, us-central1-a & Release; the fork's \evpath{.ci/dispat/gce/REPORT.md}
and the logs under \evpath{evidence/llvm-gce/logs} \\
\bottomrule
\end{longtable}
\end{center}

\subsection{RQ1: the consumer-window mechanism}
\label{sec:eval:rq1}

The consumer-relative window of \cref{def:window} is the mechanism behind
\G{2}. The mechanism campaign tests it two ways
(\evpath{campaign-mechanisms-20260910/PROTOCOL.md}). The \emph{finite
projection} enumerates every DAG consistent with one fixed order of four named
vertices, each of four single-intent sources, every completion set reachable
by recording an enabled package, and every subset of the remaining packages
designated as failed; it checks exact residual obligations, no widening, no
double ledger delivery and strict decrease exactly on progress, then applies
the source-window mutant, which discards a unit when its originating package
completes. From the \code{executor.model} object of \evpath{summary.json}: 64
graphs, 256 source instances, 772 reachable completion states, 2,270
failure-subset transitions of which 1,054 made progress, 0 projection
violations, and 260 states in which the mutant lost obligations. States and
transitions are enumeration units, not samples, and the projection is a
stable-channel, one-intent model of four vertices, not a mechanised proof.

The \emph{binary-derived counterfactual} takes each episode's original
obligation set and its observed completion tags and asks what the
source-window rule would have discarded; it runs no modified binary and
compares no competitor. \Cref{tab:eval:counterfactual} transcribes the strata.

\begin{table}[ht]
\centering\footnotesize
\caption{Source-window counterfactual by stratum: the \code{episodes},
\code{source\_mutant\_losing\_episodes} and \code{source\_mutant\_lost\_obligations}
fields of each \code{strata} entry in
\protect\evpath{experiments/saner-2027/campaign-mechanisms-20260910/summary.json}.
Units are episodes and obligations.}
\label{tab:eval:counterfactual}
\begin{tabular}{@{}llrrr@{}}
\toprule
Stratum & Value & Episodes & Losing episodes & Lost obligations \\
\midrule
Class & fault matrix & 120 & 84 & 111 \\
      & no-fault control & 10 & 0 & 0 \\
      & foreign identity & 10 & 10 & 10 \\
\midrule
Boundary & build & 40 & 28 & 41 \\
         & before publication & 40 & 28 & 35 \\
         & after artefact write & 40 & 28 & 35 \\
\midrule
Channel & stable & 70 & 47 & 61 \\
        & beta & 70 & 47 & 60 \\
\midrule
Graph & chain & 28 & 20 & 38 \\
      & diamond & 28 & 20 & 25 \\
      & fan-out & 28 & 20 & 20 \\
      & two disconnected pairs & 28 & 14 & 14 \\
      & chain with shortcuts & 28 & 20 & 24 \\
\midrule
All & & 140 & 94 & 121 \\
\bottomrule
\end{tabular}
\end{table}

\emph{Answer to \evRQ{1}.} In every state the projection reached and in every
episode the binary completed, the consumer-relative window retained exactly the
residual set; a source-relative window would have lost obligations in 260 of
the 772 states and in 84 of the 120 fault episodes (111 obligations, 70\%).
The ten no-fault controls lose nothing; the ten foreign-identity cases lose
one obligation each and test a different intervention.

\subsection{RQ2: retry and identity stability}
\label{sec:eval:rq2}

\paragraph{The mechanism matrix.} Five four-package topologies (chain,
diamond, fan-out, two disconnected pairs, chain with shortcuts), stable and
beta baselines, four fault targets and three injection boundaries give 120
fault episodes; ten no-fault controls and ten foreign-identity cases bring it
to 140. Each fault episode executes the same fault twice, repairs it and runs
a final no-op, in an order fixed by seed 20260910, with build and publish
concurrency one. The run took 113.653827~s on 9 September 2026 and passed 140
of 140 episodes (\evpath{summary.json}, \code{execution} and \code{executor}).
\Cref{tab:eval:mechanism} gives the five checks; per stratum, 40 of 40
episodes passed at each boundary, 70 of 70 in each channel and 28 of 28 on
each graph. An independently implemented offline checker reconstructed every
plan from the raw command output and re-enumerated the finite model with no
discrepancy, and a negative check on a corrupted copy failed as intended
(\evpath{RESULTS.md}). The after-write injector lets the publisher write the
artefact durably and then return a nonzero exit; it is not a crash. The
identity check is the supplied local publisher's, not a registry's.

\begin{table}[ht]
\centering\footnotesize
\caption{Mechanism checks of the 140-episode matrix, from the
\code{mechanism\_checks} object of
\protect\evpath{experiments/saner-2027/campaign-mechanisms-20260910/summary.json}.
The 130 denominator is the 120 fault episodes plus the ten foreign-identity
cases.}
\label{tab:eval:mechanism}
\begin{tabular}{@{}L{9cm}r@{}}
\toprule
Check & Passed / denominator \\
\midrule
Repeated fault preserved the residual plan and tags & 130 / 130 \\
Successful repair left an empty plan & 130 / 130 \\
Exactly one durable artefact write per obligation & 140 / 140 \\
Foreign artefact identity conflict was refused & 10 / 10 \\
Final no-op left the plan empty and tags unchanged & 140 / 140 \\
\bottomrule
\end{tabular}
\end{table}

\paragraph{Reference agreement and recovery in the companion campaigns.} The
1.10.0 campaign planned 300 generated histories, seeds 90001 to 90300, against
the restricted core oracle: 300 accepted matching rows, no rejection, invalid
execution or difference (\evpath{campaign-1.10.0-20260909/summary.json},
\code{core}); the corrected 1.9.0 campaign reports the same 300 of 300 on the
same seeds. Agreement on a fixed sample with the historical restricted model
is not a universal conformance result. The recovery stage drives 20 episodes
over graphs of 20 and 100 packages with failure probabilities .1 and .3: in
1.10.0 all 20 converged within four runs (nine episodes in two runs, five in
three, six in four) over 251 raw command records; in 1.9.0 within four runs
over 248 records (nine, six, five). The two are not pooled; the difference
is one extra attempt in the episode of size 100, probability .3, seed 100,
whose later seeded deny sets diverged, and the traces do not attribute it to
the 1.10.0 feature.

\paragraph{The 1.8.2 full-profile campaign.} On the released 1.8.2 binary, the
core shards agreed 1,000 of 1,000 and 500 of 500 with no invalid execution.
The full profile deliberately generates regressive prereleases: of 1,500
instances, 503 were refused with a validated \diag{E195} and 997 accepted; 546
of the accepted matched the historical oracle and four refused dossiers had no
row difference, so 550 comparisons were diagnostic-clean; the 1,914 divergent
package rows were classified PG1 862, PG2 375, PG2x 14 and PG3? 663
(\evpath{output/experiments-v182-summary.json}, \code{E1}). A refused
diagnostic plan is not a release plan, the oracle implements the historical
CCME 1 model, and the PG3? label is heuristic: these counts are taxonomy
against that model, not CCME 2.0.0 conformance. The first 1.8.2 attempt, whose
strict exit checker misclassified the expected refusals, is retained at
\evpath{experiments/results/v182-20260907-attempt1}; the corrected campaign
reran every seed. Appendix~B vectors 133 to 138 passed 6 of 6 on the same
binary (\code{CCME2}).

\paragraph{Historical differential, orphan, push, competitor and replay
observations.} On the released 1.4.0 binary the core profile agreed 1,000 of
1,000 and 500 of 500; the full profile's divergent rows fell into the four
predicted classes and none outside them (PG1 409, PG2 240, PG2x 58, PG3 390
on the first shard), and the 1.6.0, 1.7.0 and 1.7.2 campaigns reproduce the
per-seed rows on all 3,000 instances (\evpath{experiments/README.md},
\evpath{compare-*.md}); these validate the CCME 1.0.0 model against those
binaries, not CCME 2.0.0 admission. E2 drove a six-package npm fixture against
a Verdaccio registry behind a failing proxy: both repetitions recorded the
failed publish untagged, a same-version recovery and a final plan releasing
nothing. E5's three fresh cells passed 8 checks for the orphan scenario, 12
for a clean concurrent push and 14 for a conflicting concurrent push, on
selected interleavings rather than a sweep (\code{E2}, \code{E5}). The
competitor phases of 30 August 2026 (Lerna 10.0.1, Nx 23.1.2, Turborepo 2.10.12
with Changesets 3.0.1, release-please 17.11.2) are observations of pinned
configurations: Lerna and Nx tagged before publish, Changesets after and by
registry query, Turborepo neither versions nor publishes. E3a replayed five
releases each of two public repositories under reconstructed layouts: 36 of 42
released package-versions exact for \code{google-cloud-node} and 7 of 64 for
\code{conventional-changelog} (\evpath{experiments/results/e3a/gcn.json},
\evpath{cc.json}); neither project used the modelled intent language, and the
experiment shows the limit of policy transfer, not that those releases
followed CCME.

\emph{Answer to \evRQ{2}.} Across the designed factors of the 140-episode
matrix, recorded progress survived repeated failure without version drift,
repair emptied the plan, and the identity check refused every foreign
artefact. The restricted oracle agreed on all 300 histories of each companion
campaign and on all 1,500 core instances of the 1.8.2 campaign; the divergent
full-profile rows are retained and classified, and mark where the restricted
models stop rather than where the implementation fails.

\subsection{RQ3: availability boundaries and CI}
\label{sec:eval:rq3}

The companion campaigns execute two registry topologies and one CI fixture:
an npm package feeding a Docker base feeding a Docker application, and its
four-package extension \code{assets} (npm), \code{ui} (npm), \code{base}
(Docker), \code{app} (Docker), where \code{ui} reads its predecessor's local
\code{dist} output. Each topology repeats five configurations six times in a
frozen seeded order: mixed local-output scheduling, fully serialised controls,
a missing barrier, an injected base-publication failure with repair, and an
injected app-build failure with repair (\cref{tab:eval:rq3}).

\begin{table}[ht]
\centering\footnotesize
\caption{The 1.10.0 companion campaign, from the \code{mixed} and
\code{integration} objects of
\protect\evpath{experiments/saner-2027/campaign-1.10.0-20260909/summary.json}.
Predicates are repeated verifier checks, not samples.}
\label{tab:eval:rq3}
\begin{tabular}{@{}L{3.1cm}rrrrrr@{}}
\toprule
Topology & Executions & \shortstack[r]{Initial\\successes} &
\shortstack[r]{Designed\\failures} & Repairs & \shortstack[r]{No-op\\reruns} &
Predicates \\
\midrule
Three-package & 30 & 12 & 18 & 12 & 24 & 522 \\
Four-package & 30 & 12 & 18 & 12 & 24 & 798 \\
CI fixture & 3 repetitions & & & & & 81 \\
\bottomrule
\end{tabular}
\end{table}

Per topology the 18 designed failures are six of each negative configuration;
all 12 repairable executions completed at unchanged planned versions; the 24
completed cases performed no build or publication on their final unchanged
run. In the three-package mixed fixtures the base build overlapped npm
publication in all six observations, by 0.110716 to 0.151483~s; in the
four-package fixtures the \code{ui} build overlapped the delayed \code{assets}
publication by 0.070128 to 0.073725~s and the base build completed before the
\code{ui} publication began in all six cases, an admissible schedule
(\evpath{EXPERIMENTS-1.10.0-20260909.md}). The three CI repetitions each
cover baseline publication, a specification change, a consumer change, an
injected failing test, a missing required test, repair and an unchanged run;
both negatives blocked the release stages and the failing-test attempts wrote
no release tag. The corrected 1.9.0 campaign has the same counts and the same
522, 798 and 81 predicates. Its first attempt is preserved as failed: the
original four-node oracle demanded strict interval overlap, the traces showed
the valid schedule in which the base build completed before \code{ui}
publication began, and the predicate was corrected before the complete rerun,
whose protocol 1.10.0 then reused byte for byte (manifest SHA-256
\evhash{a3366e2fbdc590be7e8f6268bc7390146669cf7a0395837c2531a0af5554dcaa}). A
separately frozen diagnostics addendum passed 16 of 16 cases outside the
14-stage campaign.

\emph{Answer to \evRQ{3}.} In these fixtures, on one Linux ARM64 environment,
a configured barrier distinguishes local build availability from required
publication, and both a failed and an absent required test prevent release;
the campaign measures neither concurrency at scale nor onboarding effort.

\subsection{RQ4: planning cost}
\label{sec:eval:rq4}

The E4 timing cells plan generated workspaces of $P$ packages with $U$ intent
units in three graph shapes, with an optional padding history of commits that
touch no package (\cref{tab:eval:rq4}).

\begin{table}[ht]
\centering\footnotesize
\caption{Planning cost, 1.8.2 campaign: the eight rows of
\protect\evpath{experiments/results/v182-20260907/e4/timing.csv}, verbatim.
Seconds are the median, minimum and maximum of the cell's repetitions; MiB is
the maximum resident set size; two CPUs and 4~GiB on a contended
Apple-silicon Docker host.}
\label{tab:eval:rq4}
\begin{tabular}{@{}rlrrrrrrr@{}}
\toprule
$P$ & Shape & $U$ & Padding & Edges & Median s & Min s & Max s & MiB \\
\midrule
100 & layered & 200 & 0 & 185 & 0.017 & 0.017 & 0.017 & 15 \\
1,000 & layered & 200 & 0 & 2,002 & 0.250 & 0.235 & 0.252 & 21 \\
10,000 & layered & 200 & 0 & 19,919 & 19.272 & 18.992 & 20.000 & 92 \\
1,000 & chain & 200 & 0 & 999 & 0.254 & 0.241 & 0.258 & 26 \\
10,000 & chain & 200 & 0 & 9,999 & 19.354 & 18.933 & 19.645 & 118 \\
1,000 & star & 200 & 0 & 999 & 0.244 & 0.237 & 0.247 & 19 \\
1,000 & layered & 1,000 & 0 & 2,002 & 0.290 & 0.280 & 0.302 & 26 \\
1,000 & layered & 200 & 5,000 & 2,002 & 0.664 & 0.640 & 0.665 & 30 \\
\bottomrule
\end{tabular}
\end{table}

The medians rise from 0.017~s at 100 layered packages to 19.272~s at 10,000;
five-fold intent at 1,000 packages moves the median from 0.250 to 0.290~s and
5,000 padding commits to 0.664~s, the history term of the cost analysis
showing through. The earlier campaigns measured the same headline cells under
a different harness at 0.6, 2.8 to 6.7 and 110 to 149~s; those figures are not
comparable, the runner having been rewritten before the 1.8.x campaigns, and
neither set is a cross-version benchmark. The E4 recovery cells converged in
24 of 24 episodes within four runs.

\emph{Answer to \evRQ{4}.} On the capped host the observed cost grows sharply
between the 1,000- and 10,000-package cells and with history length; the
design isolates neither an asymptotic bound nor an end-to-end figure.

\subsection{RQ5: distributed execution}
\label{sec:eval:rq5}

The execution profile of \spec{28} places build and publish stages on worker
nodes that share a Git repository with the orchestrator as their only
transport. Two build graphs were measured on 22 September 2026, both on
Google Compute Engine: the Linux kernel image-build farm of
kerneltests.org~\cite{kerneltests}, whose cells are independent, and the LLVM
monorepo, whose cells consume one another's outputs. The specification's own
record of both measurements is the entry ``2026-09-22: Distributed execution
implemented'' of \code{DESIGN-HISTORY.md}, which this subsection follows. The
claim the measurements support is narrow and is stated once: on these
fixtures, once each, a distributed release took no longer than the same
recipe on one node; no speedup of the profile or of the engine is claimed
beyond that. The raw observations are copied unedited into
\evpath{experiments/distributed-20260922/evidence/}, whose
\evpath{MANIFEST.json} records every file's origin and SHA-256 and whose
\evpath{CROSS-CHECK.json} traces every headline figure below to a raw file
(43 checks, 41 comparable, all agreeing). Each fork carries its own acceptance
script: the kernel fork's \evpath{dispat/compare.py} accepts a run only if its
cases, revision, compiler and assembler identities match the public
reference, every cell succeeded, and the elapsed time is no greater than the
reference; the LLVM fork's \evpath{.ci/dispat/gce/compare.py} compares the
install trees file by file with the prefix masked.

\subsubsection{The kernel farm}
\label{sec:eval:rq5:kernel}

The public reference for each group is the build step of the upstream builder
for kernel v7.3-rc4 at \evhash{93f51579e7df248780214094418f205253383cc5}:
alpha 623~s, hexagon 111~s and m68k\_nommu 128~s, with GCC 16.2.0 over GNU
assembler 2.47.20260726 from the kernel.org crosstool
builds~\cite{crosstool} and Ubuntu clang 18.1.3 for hexagon
(\evpath{evidence/kernel-gce/round1-codex/baseline.json}). These are
build-step durations on the farm's own, different, hardware. The farm's
configuration gives image builds an exclusive lock, one per machine
(\evpath{master/master.cfg}, unchanged in the fork from upstream
\code{336111e3}), and its maintainer reports that the build directory is a RAM
disk that holds a single build (correspondence of 23 September 2026). The
fork's first report describes the reference steps as shared with another
build; the configuration does not support that. The references bound
acceptance and enter no other comparison. A first run on 22 September established the coordination profile
and two engine defects and is not parity evidence (wrong compiler, results
not retained). Round 1 then installed the reference toolchains on a
\code{c3-standard-22} and a \code{t2d-standard-8} worker and measured a
grouped control and four candidates, each once; every candidate failed the
gate on elapsed time alone (split hexagon 170.07~s, split m68k\_nommu
159.03~s, batched hexagon 165.30~s, batched m68k\_nommu 196.23~s; the control
1,797.75~s with alpha's script alone at 1,324~s against 623), with cases,
toolchains and revision accepted (\cref{tab:eval:kernel1}; the curated archive
of 145 files, SHA-256 \code{2fcc02ec\ldots}, is under
\evpath{evidence/kernel-gce/round1-codex/}).

Round 2 calibrated the farm's script bodies outside dispat on prepared source
trees (\evpath{round2/calibration/cal-c3.txt}): the three hexagon
configurations started together at 11 jobs each finished after 120.12~s, the
five m68k\_nommu configurations at 7 jobs each after 128.17~s, both at or above
their references before any release overhead. It then pinned each group's
costliest cells to the \code{c3-standard-22} that starts the release and gave
the rest to the \code{t2d-standard-8} worker, set \code{-j} per cell from the calibration instead of the
script's default of 1.5 times the processor count, leased a prepared source
tree per build slot instead of a per-cell worktree, and recorded the worker's
cell on the orchestrator. The build commands are the farm's; the fork's
\code{stable-build-arch.sh} adds only the environment overrides for the build
folder, the log, the job count and the compiler directory, and the selection
of a subset of configurations. \Cref{tab:eval:kernel2} transcribes every timed release
of the round, including the rejected first candidate and a rerun that planned
nothing.

\begin{table}[ht]
\centering\footnotesize
\caption{Kernel round 2, every timed release. Elapsed is the GNU \code{time}
record (\protect\evpath{par-c3/*.time}, \protect\evpath{solo-c3/*.time}, also
\protect\evpath{*.elapsed}); script and lease seconds per cell and the
\code{compare.py} verdict with its ratio (reference over elapsed) are from
\protect\evpath{verdicts.txt}; all under \protect\evpath{dispat/gce/round2} of
the fork and under
\protect\evpath{experiments/distributed-20260922/evidence/kernel-gce/fork-codex-kernel-parity-c699bf3d/dispat/gce/round2}.
Each row is one release.}
\label{tab:eval:kernel2}
\begin{tabular}{@{}L{2.4cm}L{4.9cm}rrL{3.3cm}@{}}
\toprule
Run & Cells: script s (lease s) & Elapsed s & Ref.\ s & Verdict, ratio \\
\midrule
hexagon, first candidate & c3: tinyconfig 60 (12.0), defconfig 84 (12.0);
w1: allnoconfig 80 (12.0) & 111.47 & 111 & reject, elapsed only; 0.9958 \\
hexagon, rerun of the unchanged configuration & no cell built: every package
reported \code{unchanged} & 3.83 & 111 & not compared; planned nothing \\
hexagon, trial 1 & c3: tinyconfig 70 (0.0), defconfig 85 (0.0); w1: allnoconfig
82 (0.0) & 95.89 & 111 & accept; 1.1576 \\
hexagon, trial 2 & c3: tinyconfig 71 (0.0), defconfig 84 (0.0); w1: allnoconfig
81 (0.0) & 92.22 & 111 & accept; 1.2036 \\
m68k\_nommu, trial 1 & c3: m5249evb 95 (1.0), m5272c3 96 (1.0), m5407c3 96
(14.0), m5307c3 97 (15.0); w1: m5475evb 50 (0.0) & 115.61 & 128 & accept;
1.1072 \\
m68k\_nommu, trial 2 & c3: m5249evb 105 (0.0), m5407c3 105 (0.0), m5307c3 106
(0.0), m5272c3 107 (0.0); w1: m5475evb 51 (0.0) & 115.48 & 128 & accept;
1.1084 \\
hexagon, one node alone & solo: tinyconfig 91 (13.0), allnoconfig 105 (13.0),
defconfig 110 (12.0) & 122.62 & 111 & reject, elapsed only; 0.9052 \\
m68k\_nommu, one node alone & solo: m5475evb 56 (0.0), m5407c3 91 (0.0),
m5307c3 92 (0.0), m5272c3 73 (58.0), m5249evb 72 (66.0) & 138.01 & 128 &
reject, elapsed only; 0.9275 \\
\bottomrule
\end{tabular}
\end{table}

\code{compare.py} accepted the exact case lines, compiler, assembler and
revision in all seven releases that built cells and rejected three on elapsed
time alone; every cell reported \code{passed}. The gate was met twice for
hexagon after one rejected candidate at 111.47~s (a 12.0~s lease wait on each
of its three cells, before the leased-tree change), and twice for m68k\_nommu
with none. The same-hardware pairs are the single-node rows: each group run
once more on one \code{c3-standard-22} with no worker and nothing else
changed, rejected on elapsed time alone at 122.62 and 138.01~s. Two features
of those runs are stated rather than hidden: a release that names no worker
is not bounded by \code{execution.concurrency} and pays no preflight, so its
figure is the best the machine can do with the recipe; and the lease check
dominated the m68k\_nommu run, two cells recording 58 and 66~s of
\code{source\_lease\_seconds} against build scripts of 73 and 72~s
(\evpath{solo-c3/results/*/result.json}), so 138.01~s is pessimistic by some
part of that cost. The calibration outside dispat is the corroboration. alpha
was not run: its \code{allmodconfig} is one cell that cannot be split, with a
lower bound of about 1,198~s against a budget of about 610~s on the largest
machine the project's 32-vCPU quota allows; a group the experiment could not
measure, not one it measured and failed. The engine has two placement classes
and no per-worker pinning, which sufficed because each group needed exactly
two machines; the worker sat idle 59 and 49~s in the m68k\_nommu trials while
the orchestrator finished four cells. Fixed overhead per release was 3 to 8~s
of preflight; the round's five VMs lived 10 to 22 minutes each.

\paragraph{Why these layouts do not transfer to the farm.} The paired
layouts ran two hexagon or four m68k\_nommu image builds at once on the
\code{c3-standard-22}, the single-node runs three and five, and the
calibration the same stacking, each build in its own directory under a
\code{/tmp} on the boot disk (the image's default; no mount output was
retained). The farm runs one image build per machine in a RAM disk that holds
one build, so none of these layouts fits its machines as they are, and the
measured build directories were disk-backed where the farm's are in memory.
The round shows what a second machine bought on this hardware; it does not
show that the farm would finish sooner.

\subsubsection{The LLVM pair}
\label{sec:eval:rq5:llvm}

The fork \code{yohimik/llvm-project}, branch \code{dispat-gce}, LLVM
\code{main} at \code{a9f4d3ff}, builds with \code{LLVM\_FAST=1} (assertions,
examples, docs and benchmarks off), \code{LLVM\_SKIP\_CHECKS=1}, the X86
target, 14 build jobs and 2 link jobs, on an \code{e2-standard-4} orchestrator
and two \code{n2-custom-14-49152} workers (7 cores each), Ubuntu 24.04.5,
g++ 13.3.0, cmake 3.28.3, ninja 1.11.1. The recipe's identity record embeds
the toolchain versions, the source revision and its own hash into every
receipt, and a consumer refuses a provider whose receipt does not match the
installed tree. Nine attempts were made; six have a retained orchestrator log
and five a GNU \code{time} record (\evpath{evidence/llvm-gce/logs/}, mapped
to attempts in \evpath{MANIFEST.json}; \cref{tab:eval:llvm-attempts}). Six
were discarded or failed for reasons recorded with them: the recipe's
fingerprint hashed umask-dependent modes; an orchestrator without a C++
compiler; a baseline clone without the \code{v24.0.0} tags; a fixture change;
and a dispat defect, captured outputs passed through the checkout's
\code{.gitattributes} text conversion, which rewrote CRLF pairs inside three
libraries and failed the bundle in attempts 4 and 7, fixed on \code{main} at
\code{9f649f88}. A second engine defect, a fixed 30~s report bound that killed
the 1.4~GB result push, was fixed at \code{f0c9e959} before the pair. The
accepted pair is attempts 8 and 9.

The single-heavy fixture of attempts 4 and 5 (\code{llvm}, then \code{lld}
and \code{polly} consuming it) is the negative case \spec{28.7} predicts: on
one node the three were published at 20:16 (whole release 36:09.09 with a
15:53 xz bundle); distributed, at 26:47.69, a wall time that ends in exit 1 on
the bundle. One package holding nine tenths of the work cannot pay for moving
its 1.4~GB install tree to consumers that build in two minutes. The matrix
fixture adds \code{llvm-aarch64}, the same sources built for AArch64 as a
package that depends on nothing, so that one node must serialise two heavy
builds and a pool need not, and sets an uncompressed bundle on both sides
(\cref{tab:eval:llvm-matrix}).

\begin{table}[ht]
\centering\footnotesize
\caption{The LLVM matrix pair, attempts 8 and 9. Per-package figures are the
results table of \protect\evpath{.ci/dispat/gce/REPORT.md}; the whole-release
wall times are the \code{Elapsed (wall clock) time} lines of
\protect\evpath{release-local3.time} and \protect\evpath{release-dist6.time}
under \protect\evpath{experiments/distributed-20260922/evidence/llvm-gce/logs/},
whose \code{done} lines record 5 published and 0 failed on both sides. One
repetition per side.}
\label{tab:eval:llvm-matrix}
\begin{tabular}{@{}L{4.6cm}L{3.6cm}L{5cm}@{}}
\toprule
& Baseline, \code{w2} alone & Distributed, \code{w1} + \code{w2} +
\code{lb-orch} \\
\midrule
\code{llvm} (X86) & 38:38, sharing the node & 22:08 on \code{w1}, including
the 1.4~GB result push \\
\code{llvm-aarch64} & 38:51, sharing the node & 23:22 on \code{w2}, including
its result push \\
\code{lld} + \code{polly} & 2:16 & 2:13 on the two workers, after \code{llvm}'s
outputs reached them \\
bundle (uncompressed tar) & 0:18 & 0:32 on \code{lb-orch} \\
whole release, GNU \code{time} & 41:12.02, exit 0 & 27:48.94, exit 0 \\
\bottomrule
\end{tabular}
\end{table}

\code{llvm}'s declared outputs, 1.4~GB in 2,647 files, were pushed once and
fetched by two consumers and the orchestrator; \code{lld} (420~MB) and
\code{polly} (22~MB) consumed them on nodes that never built \code{llvm},
which is the mechanism \spec{28.7} names as the profile's required gain. The
bundle verified every provider's receipt against the installed tree, so the
transport is byte-faithful; \code{compare.py} then finds the two sides differ
only in what the recipe embeds about its checkout (the compiled-in VCS
revision and the build-id derived from it, the repository string, RPATH and
\code{\_\_FILE\_\_} paths), and \code{libLLVMCore.a}, which embeds none of
these, is byte-identical; byte identity throughout needs
\code{LLVM\_APPEND\_VC\_REV} off, a file prefix map and a relative RPATH,
which this recipe does not set. Coordination was cheap beside the builds: probe
1 to 4~s, claim 3 to 6~s warm, a result noticed within 5~s; a cold worker's
first claim comes minutes later because its poll fetches the snapshot branch
first, about 10~min for the 3.6~GB history.

\paragraph{Binaries and caveats.} Attempt 9 ran
\code{1.11.0-rc.4+main.272772cb}, built in a checkout with uncommitted
documentation edits, so that Go stamped it \code{vcs.modified=true}; the
deployed binary is recorded by its SHA-256 prefix \code{feba0210d13c90fa},
and the reproduction kit's clean rebuild of the same commit gives
\evhash{46bc2cf821dd57479c27bf2a4b7eee63b69e01db6f18fbd6e663ad66e6812ba8},
differing only by that stamp; \code{f0c9e959} rebuilds byte-identically to
\evhash{63441cc66f68d07b5efbf9a42edd5559304d153d48b82b1dfb037349e20606ed}.
Both measured builds credited a unit's whole source set to every reached
dependent until product \code{main} \code{e461d7d9}; the fixtures carry
single-source units only, which both rules attribute identically
(\cref{ob:attribution}). The qualifications the specification's own entry
states are repeated without softening: one repetition per side on both
graphs, so variability is unmeasured; the kernel's second node is a slower
machine of another family, so each kernel pair is one node against that node
plus a second, slower one; the kernel groups have only independent cells to
distribute; the LLVM matrix's two heavy builds are one project for two
targets; the LLVM distributed side had two build nodes against one; and the
single-node m68k\_nommu figure includes 58 to 66~s of lease contention. The
entry dates the LLVM pair 2026-09-23 while every retained log is stamped
2026-09-22 UTC, and names the orchestrator-side commit \code{6c1a5fcd} where
the fork's is \code{d8fb8408}; both are recorded as observed, the driver
record \evpath{llvm-gce/drivers/pair4.out} showing the second to be the same
change applied on the orchestrator as a patch.

\emph{Answer to \evRQ{5}.} On the kernel fixture the distributed release met
the public reference gate twice for each of two groups (hexagon 95.89 and
92.22~s against 111, after one rejected candidate; m68k\_nommu 115.61 and
115.48~s against 128) where the same recipe on the same single node did not
(122.62 and 138.01~s), each single-node run once; alpha was outside the
project's capacity. On the LLVM fixture the distributed release of a
five-package matrix took 27:48.94 against 41:12.02 on one node, produced the
same five products under the recipe's receipts, and the negative case
predicted by \spec{28.7} was observed at 26:47.69 against 20:16. On these
fixtures, once each, distributed took no longer than single node; nothing
more is claimed.

\subsection{Threats to validity}
\label{sec:eval:threats}

The mechanism matrix is a designed, deterministic set of 140 episodes with no
population inference; its after-write fault is an adapter response, not a
crash; the projection is four vertices, one intent and the stable channel;
the counterfactual runs no modified binary; the runner image cannot be
rebuilt bit for bit on a fresh machine because \code{apk add} consults an
external repository. The companion campaigns' oracle is the restricted core
profile; their predicate counts are not independent observations; their
delays expose availability, not throughput; their environment is one Linux
ARM64 Docker host and validates nothing about native Windows, which no
campaign measures. The E1 oracle is the CCME 1.0.0 model, so its divergence
classes are taxonomy against that model; the E3a replay reconstructs layouts
the projects never declared; planning-cost figures are host-sensitive and
no cross-version benchmark. The kernel rounds are one pair of machine shapes,
one region, one kernel revision, two trials per accepted topology and one
single-node run per group, gated against public build-step times from other
hardware, with image builds stacked on one machine in disk-backed directories
where the farm runs one per machine in a RAM disk; alpha, every QEMU group and the 57-group matrix are unmeasured. The
LLVM pair is one repetition per side with two build nodes against one, the
same sources for two targets, a recipe without the check targets, six of nine
attempts failed or discarded, a deployed binary recorded by hash prefix only,
and raw logs copied from the measuring session's scratchpad into the evidence
tree rather than kept in the fork.

\subsection{Reproduction}
\label{sec:eval:repro}

The Docker and Terraform kit for the distributed measurements is at
\evpath{experiments/distributed-20260922/} (README, \evpath{pins.json}, node
images, Terraform roots, runbooks, a local Compose check), beside the
\evpath{evidence/} directory. Verified on 22 September 2026 on an Apple
silicon workstation: the experiment binaries rebuild byte for byte from their
commits in a digest-pinned \code{golang:1.26.5-alpine} to the hashes above;
the four round-2 configurations are regenerated by the fork's generator byte
for byte apart from the mailbox path; both Terraform roots validate and plan
against the original projects (thirteen resources each) but \code{terraform
apply} was not run, so no cloud pair was re-measured; and the pipeline ran end
to end on the workstation with two containers of the kernel node image, the
fork's \code{compare.py} accepting the run (three cases, revision
\code{93f51579}, 63.47~s, a laptop figure that says nothing about parity;
\evpath{runs/local-kernel-20260922T184954Z/verdict.json}). A reproducibility
audit of every older campaign, with rerun commands and status per family, is
\evpath{experiments/REPRODUCIBILITY.md}.
